\documentclass[12pt]{article}
\usepackage{amsmath,amssymb,amsthm}
\usepackage{booktabs}
\usepackage[margin=1in]{geometry}

\newtheorem{theorem}{Theorem}
\newtheorem{lemma}[theorem]{Lemma}
\newtheorem{corollary}[theorem]{Corollary}

\title{Problem 448: Average Least Common Multiple}
\author{}
\date{}

\begin{document}
\maketitle

\section{Problem Statement}
Define $f(n) = \frac{1}{n}\sum_{k=1}^{n} \operatorname{lcm}(k, n)$. Compute $S(N) = \sum_{n=1}^{N} f(n) \pmod{10^9 + 7}$ for $N = 10^8$.

\section{Mathematical Foundation}

\begin{theorem}[Closed Form]
\label{thm:closed}
For all $n \geq 1$,
\[
f(n) = \frac{1 + \displaystyle\sum_{d \mid n} d\,\varphi(d)}{2}.
\]
\end{theorem}

\begin{proof}
Using $\operatorname{lcm}(k,n) = kn/\gcd(k,n)$:
\[
f(n) = \sum_{k=1}^{n} \frac{k}{\gcd(k,n)}.
\]
Group by $g = \gcd(k,n)$ with $k = gj$, $\gcd(j, d) = 1$, $d = n/g$:
\[
f(n) = \sum_{d \mid n} \Phi(d), \quad \Phi(d) := \sum_{\substack{j=1\\\gcd(j,d)=1}}^{d} j.
\]
For $d \geq 2$, the coprime residues pair as $j \leftrightarrow d - j$ (both coprime to~$d$, with $j \neq d-j$ for $d \geq 3$; for $d = 2$, $\Phi(2) = 1$). Each pair sums to~$d$, and there are $\varphi(d)/2$ pairs, so $\Phi(d) = d\,\varphi(d)/2$. Since $\Phi(1) = 1 = 1 \cdot \varphi(1)/2 + 1/2$:
\[
f(n) = \frac{1 + \sum_{d \mid n} d\,\varphi(d)}{2}.\qedhere
\]
\end{proof}

\begin{corollary}[Summation Formula]
\[
S(N) = \frac{N + \displaystyle\sum_{d=1}^{N} d\,\varphi(d)\,\lfloor N/d \rfloor}{2}.
\]
\end{corollary}

\begin{proof}
Sum Theorem~\ref{thm:closed} over $n = 1, \ldots, N$ and swap the order of summation:
\[
\sum_{n=1}^{N}\sum_{d \mid n} d\,\varphi(d) = \sum_{d=1}^{N} d\,\varphi(d)\,\lfloor N/d \rfloor.\qedhere
\]
\end{proof}

\begin{lemma}[Multiplicativity]
$h(n) = n\,\varphi(n)$ is multiplicative with $h(p^a) = p^{2a} - p^{2a-1}$.
\end{lemma}

\begin{proof}
$\operatorname{Id}$ and~$\varphi$ are both multiplicative, so their product is multiplicative. For prime powers: $h(p^a) = p^a \cdot p^{a-1}(p-1) = p^{2a-1}(p-1)$.
\end{proof}

\begin{lemma}[Coprime Residue Pairing]
For $d \geq 2$ and $\gcd(j,d) = 1$ with $1 \leq j \leq d$: $\gcd(d-j, d) = 1$ and $j + (d-j) = d$, with $j \neq d - j$ for $d \geq 3$.
\end{lemma}

\begin{proof}
$\gcd(d-j, d) = \gcd(j, d) = 1$. For $j = d/2$: $\gcd(d/2, d) = d/2 = 1$ requires $d = 2$.
\end{proof}

\section{Concrete Values}

\begin{center}
\begin{tabular}{cccc}
\toprule
$n$ & $\sum_{d \mid n} d\varphi(d)$ & $f(n)$ & $S(n)$ \\
\midrule
1 & 1  & 1  & 1 \\
2 & 3  & 2  & 3 \\
3 & 7  & 4  & 7 \\
4 & 11 & 6  & 13 \\
5 & 21 & 11 & 24 \\
6 & 21 & 11 & 35 \\
\bottomrule
\end{tabular}
\end{center}

\section{Algorithm}

\begin{enumerate}
\item \textbf{Totient sieve:} Compute $\varphi(n)$ for $n = 1, \ldots, N$ in $O(N \log\log N)$.
\item \textbf{Prefix sums:} Compute $H(k) = \sum_{d=1}^{k} d\,\varphi(d) \bmod p$ for $k = 1, \ldots, N$.
\item \textbf{Grouped summation:} Evaluate $\sum_{d=1}^{N} d\,\varphi(d)\,\lfloor N/d \rfloor$ by grouping the $O(\sqrt{N})$ distinct values of $\lfloor N/d \rfloor$.
\item \textbf{Final:} $S(N) = (N + \text{sum}) \cdot 2^{-1} \bmod p$.
\end{enumerate}

\section{Complexity Analysis}

\begin{itemize}
\item \textbf{Time:} $O(N)$ for the totient sieve and prefix sums; $O(\sqrt{N})$ for the grouped summation. Total: $O(N)$.
\item \textbf{Space:} $O(N)$ for the $\varphi$ and prefix-sum arrays.
\end{itemize}

\section{Answer}

\[
\boxed{106467648}
\]

\end{document}
